\section{Next Greater Element II (Circular Array)}
\label{sec:sq-next-greater-2}

\problemstatement{We are given a \textbf{circular} array $\textit{nums}$
(the element after the last one wraps around to the first). For every
element, find its next greater element, searching circularly -- i.e.\ we
may wrap around the end of the array once, but not more than once.}

\begin{patternbox}
The keyword \emph{``circular''} attached to a next-greater-element problem
signals that the search for each element's answer may need to continue
past the array's literal end and wrap to the beginning -- but only
\textbf{once around}, not indefinitely. The standard trick for simulating
one extra lap around a circular array without actually duplicating it in
memory is to iterate over \textbf{twice} the indices ($0$ to $2n-1$), using
the modulo operation to map each iteration index back into the real array.
\end{patternbox}

\subsection*{Understanding the problem carefully}

If we simply ran the Section~\ref{sec:sq-next-greater-1} algorithm once
over $\textit{nums}$, an element near the \emph{end} of the array would
never get a chance to find a next-greater candidate that happens to sit
near the \emph{beginning} -- the algorithm would incorrectly report $-1$
for it. Conceptually, we want to pretend the array is duplicated end to
end ($\textit{nums} + \textit{nums}$) and search within that doubled view,
but stop considering circular wraparound after one full lap (we must not
wrap twice, or an element could incorrectly become its own next greater
element).

\begin{ideabox}[Simulate a Doubled Array via Modulo Indexing]
Run the same right-to-left monotonic-stack scan as
Section~\ref{sec:sq-next-greater-1}, but let $i$ range from $2n-1$ down
to $0$, accessing $\textit{nums}[i \bmod n]$ at each step. Only
\emph{record} an answer for $i < n$ (the first lap, which corresponds to
real array positions); the second lap ($i \ge n$, processed first since we
scan from $2n-1$ downward) exists purely to \emph{seed the stack} with
candidates that wrap around from the beginning of the array, without
double-counting any element as its own answer.
\end{ideabox}

\subsection*{Why processing the second lap first, without recording answers, is correct}

Scanning from $2n-1$ down to $0$ means we process the conceptual ``second
copy'' of the array (indices $n$ to $2n-1$, mapping to real indices $0$
to $n-1$) \emph{before} the conceptual ``first copy'' (indices $0$ to
$n-1$). By the time we reach real index $i$ in the first-copy pass (i.e.\
when the loop variable itself is also $i$, having already processed
$i+1, \dots, 2n-1$), the stack already reflects every element that could
serve as $i$'s next greater element either later in the same lap
\emph{or} wrapped around from the start of the array (since those wrapped
candidates were pushed during the second-lap pass that ran first). This
correctly captures one full circular lookahead without ever risking a
second wraparound, since each real index contributes to the stack at most
twice (once per lap) and answers are only recorded during the final,
first-copy pass.

\subsection*{Algorithm}

\begin{enumerate}
  \item Initialize an empty stack and a result array of size $n$, filled
        with $-1$.
  \item For $i$ from $2n-1$ down to $0$:
  \begin{enumerate}
    \item Let $\textit{val} \gets \textit{nums}[i \bmod n]$.
    \item While the stack is not empty and its top $\le \textit{val}$:
          pop.
    \item If $i < n$: set $\textit{result}[i] \gets$ the stack's top if
          non-empty, else $-1$.
    \item Push \textit{val}.
  \end{enumerate}
  \item Return \textit{result}.
\end{enumerate}

\begin{algorithm}[H]
\caption{Next Greater Element II}
\begin{algorithmic}[1]
\Procedure{NextGreaterElementsCircular}{\textit{nums}}
  \State $\textit{result}[0..n-1] \gets -1$, $\textit{stack} \gets$ empty
  \For{$i \gets 2n-1$ \textbf{down to} $0$}
    \State $\textit{val} \gets \textit{nums}[i \bmod n]$
    \While{\textit{stack} not empty \textbf{and} top $\le \textit{val}$}
      \State pop
    \EndWhile
    \If{$i < n$ \textbf{and} \textit{stack} not empty}
      \State $\textit{result}[i] \gets$ top
    \EndIf
    \State push \textit{val}
  \EndFor
  \State \Return \textit{result}
\EndProcedure
\end{algorithmic}
\end{algorithm}

\subsection*{Dry run}

\dryrun{Take $\textit{nums} = [1,2,1]$ ($n=3$).}

\begin{itemize}
  \item $i=5$ ($\textit{val}=\textit{nums}[2]=1$): stack empty. $i\ge n$,
        no record. Push $1$. Stack $=[1]$.
  \item $i=4$ ($\textit{val}=\textit{nums}[1]=2$): top $1\le2$, pop.
        Stack empty. $i\ge n$, no record. Push $2$. Stack $=[2]$.
  \item $i=3$ ($\textit{val}=\textit{nums}[0]=1$): top $2>1$, no pop.
        $i\ge n$, no record. Push $1$. Stack $=[2,1]$.
  \item $i=2$ ($\textit{val}=\textit{nums}[2]=1$): top $1\le1$, pop.
        New top $2>1$, no further pop. $i<n$: $\textit{result}[2]=2$.
        Push $1$. Stack $=[2,1]$.
  \item $i=1$ ($\textit{val}=\textit{nums}[1]=2$): top $1\le2$, pop. New
        top $2\le2$, pop. Stack empty. $i<n$:
        $\textit{result}[1]=-1$. Push $2$. Stack $=[2]$.
  \item $i=0$ ($\textit{val}=\textit{nums}[0]=1$): top $2>1$, no pop.
        $i<n$: $\textit{result}[0]=2$. Push $1$. Stack $=[2,1]$.
\end{itemize}

Result: $[2,-1,2]$. Checking by hand: index $0$ (value $1$) wraps to find
$2$ at index $1$; index $1$ (value $2$) has no greater element anywhere
(it is the maximum); index $2$ (value $1$) wraps around to find $2$ at
index $1$ (or directly via the doubled view, finds the $2$ that appears
right after it when wrapping).

\subsection*{C++ implementation}

\begin{lstlisting}
#include <bits/stdc++.h>
using namespace std;

vector<int> nextGreaterElementsCircular(vector<int>& nums) {
    int n = nums.size();
    vector<int> result(n, -1);
    stack<int> st;

    for (int i = 2 * n - 1; i >= 0; i--) {
        int val = nums[i % n];
        while (!st.empty() && st.top() <= val) {
            st.pop();
        }
        if (i < n && !st.empty()) {
            result[i] = st.top();
        }
        st.push(val);
    }
    return result;
}

int main() {
    vector<int> nums = {1, 2, 1};
    for (int x : nextGreaterElementsCircular(nums)) cout << x << " ";
    cout << endl;
    return 0;
}
\end{lstlisting}

\begin{complexitybox}
\textbf{Time Complexity.} The loop runs $2n$ iterations, each doing
amortized $O(1)$ stack work (every element is pushed at most twice and
popped at most twice across the whole run), giving
\[
  O(n).
\]
\textbf{Space Complexity.} $O(n)$ for the stack and result array.
\end{complexitybox}

\begin{takeawaybox}
``Circular array, search wrapping around once'' has a standard, reusable
fix: iterate over $2n$ virtual indices with modulo wraparound, using the
first lap purely to seed state (here, the monotonic stack) and the second
lap to record real answers. This avoids physically duplicating the array
and generalizes to other circular-array problems beyond monotonic stacks.
\end{takeawaybox}
